\section{Technische Anforderungen}
Die technischen Anforderungen beschreiben die notwendigen Eigenschaften und Randbedingungen für die Umsetzung des Systems.
Sie werden in die Bereiche Hardware, Software, Schnittstellen und Performance unterteilt.

\subsection{Hardware}
Die Hardwarearchitektur basiert auf einem Mikrocontroller mit integrierter \gls{wlan}-Schnittstelle.
Dieser übernimmt sowohl die Datenverarbeitung als auch die Kommunikation mit externen Diensten.
Die Anzeige erfolgt über mehrere \gls{led}-Stränge die aus vorgefertigen \gls{led}-Streifen von AZ-Delivery bestehen und mit \gls{ws2812b} \gls{rgb}-Leuchtdiode bestückt sind~\cite{AZDelivery}.
Ein \gls{led}-Strang wird für die Hauptanzeige verwendet. Das ist die Uhrzeit in Stunden und Minuten, die Außentemperatur und das Wetter.
Der zweite \gls{led}-Strang wird für die Sekundenanzeige verwendet.
Der dritte \gls{led}-Strang wird für die Anzeige der Tageszeit, der Raumfeuchtigkeit, der Mondphase und der Innenraumtemperatur.
Alle Hardware-Anforderungen sind in Tabelle~\ref{tab:Hardware Anforderung} aufgelistet.

\begin{table}[H]
\caption{Technische Anforderungen – Hardware}
\centering
\begin{tabular}{p{0.5cm} p{12.5cm}}
\textbf{ID} & \textbf{Beschreibung} \\
\hline
H1 & Verwendung eines Mikrocontrollers mit WLAN-Funktionalität \\
H2 & Ansteuerung von einzeln adressierbaren LEDs \\
H3 & 3 Kaskadierte LED Strukturen mit je einem Datenpin am Mikrocontroller \\
H4 & Einbindung eines Temperatur- und Feuchtigkeitssensors \\
H5 & Sicherstellung einer stabilen Spannungsversorgung für alle Komponenten \\
\end{tabular}
\label{tab:Hardware Anforderung}
\end{table}

\subsection{Software}
Die Softwarearchitektur ist modular aufgebaut.
Zentrale Komponenten sind die Zeitverarbeitung, die \gls{led}-Ansteuerung sowie die Kommunikation mit externen Diensten.
Die Verarbeitung erfolgt regelbasiert und ermöglicht eine flexible Erweiterung.
Alle Software-Anforderungen sind in Tabelle~\ref{tab:Software Anforderung} aufgelistet.

\begin{table}[H]
\caption{Technische Anforderungen – Software}
\centering
\begin{tabular}{p{0.5cm} p{12cm}}
\textbf{ID} & \textbf{Beschreibung} \\
\hline
S1 & Implementierung einer Zeitverwaltung mit NTP-Synchronisation \\
S2 & Regelbasierte Umwandlung numerischer Zeitwerte in Wortdarstellung \\
S3 & Verarbeitung und Klassifikation von Wetterdaten \\
S4 & Implementierung einer Sensorabfrage für Innenraumwerte \\
S5 & Steuerung der LED-Anzeige in Echtzeit \\
S6 & Bereitstellung eines HTTP-Servers zur Benutzerinteraktion \\
S7 & Fehlerbehandlung bei Kommunikations- und Sensorfehlern \\
\end{tabular}
\label{tab:Software Anforderung}
\end{table}

\subsection{Schnittstellen}
Die Schnittstellen ermöglichen die Integration externer Datenquellen sowie die Interaktion mit dem System.
Die Kommunikation erfolgt überwiegend über standardisierte Protokolle.
Alle Schnittstellen-Anforderungen sind in Tabelle~\ref{tab:Schnittstellen Anforderung} aufgelistet.

\begin{table}[H]
\caption{Technische Anforderungen – Schnittstellen}
\centering
\begin{tabular}{p{0.5cm} p{12.5cm}}
\textbf{ID} & \textbf{Beschreibung} \\
\hline
I1 & WLAN-Schnittstelle zur Internetanbindung \\
I2 & HTTP-Schnittstelle zur Benutzerinteraktion per \gls{ui} \\
I3 & API-Anbindung zur Abfrage von Wetterdaten \\
I4 & Sensorschnittstelle zur Erfassung von Temperatur und Luftfeuchtigkeit \\
\end{tabular}
\label{tab:Schnittstellen Anforderung}
\end{table}

\subsection{Performance}
Die Anforderungen an die Performance stellen sicher, dass das System eine aktuelle und konstante Anzeige gewährleistet.
Gleichzeitig werden externe Daten in definierten Intervallen aktualisiert, um eine ausreichende Aktualität zu erreichen.
Alle Performance-Anforderungen sind in Tabelle~\ref{tab:Performance Anforderung} aufgelistet.

\begin{table}[H]
\caption{Technische Anforderungen – Performance}
\centering
\begin{tabular}{p{0.5cm} p{10.5cm}}
\textbf{ID} & \textbf{Beschreibung} \\
\hline
P1 & Aktualisierung der Anzeige in Intervallen kleiner gleich 1 s \\
P2 & Wetterdatenaktualisierung mindestens alle 15 min \\
P3 & Sensoraktualisierung im Bereich weniger Sekunden \\
P4 & Minimierung von Verzögerungen bei der Benutzerinteraktion \\
\end{tabular}
\label{tab:Performance Anforderung}
\end{table}